%\DeclareOption{s5paper}
%    {\setlength\paperheight {240mm}%
%     \setlength\paperwidth  {165mm}}
% S5: 165 x 242 mm (This is also called statsformat (in Swedish), it.s a special government format used in Sweden) 
% Standard format for master thesis works, according to LIU-tryck

\documentclass[]{report}
%\documentclass[a4paper]{report}
%\documentclass[s5paper]{report}

% Available sectioning commands for report: \part{}, \part{}, \chapter{}, \section{}, \subsection{}, \paragraph{}, \subparagraph{}.

% "-": Hyphen, used to join words and to separate syllables in words
% "$-$": Minus
% "--": En-dash, used for number ranges, to mark off nested clausees and phrases (use)
% "---": Em-dash, used to mark off nested clausees and phrases (don't use)

% % % % % % %
% PREAMBLE  %
% % % % % % %

% Packages

% Miscellaneous tools
\usepackage{etoolbox} % Includes the macro \ifstrequal
% Paper size
\iffalse
    \usepackage{./Packages/s5paper/s5paper}
\fi
% Linking
\iftrue
    \usepackage{hyperref} % Links to pages
    \hypersetup{pdfborder = {0 0 0}} % Remove the frame around links
\else
    \usepackage{url} % Doesn't link to pages
\fi
% Graphics
\usepackage{graphicx} % Necessary in order to include images
\usepackage{titlepic} % Necessary in order to be able to have an image on the title page
% Referencing
\iftrue
    \usepackage[round, authoryear]{natbib} % Harward style referencing
\else
    \usepackage[square, numbers]{natbib} %
\fi

%Define commands
\newcommand{\comment}[1]{}
\newcommand{\HRule}{\rule{\linewidth}{0.5mm}}

\begin{document}

% % % % % % % %
% TOP MATTER  %
% % % % % % % %

%\begin{titlepage}
% Top matter
\title{Simulating Ocean Waves and Ships on a Tree Structure}
\author{Kristofer Krus}
\date{\today}
\titlepic{\includegraphics[width=0.6\textwidth, clip=true]{./Images/lith_en_vert_col.pdf}}
%\includegraphics[width=0.6\textwidth, clip=true]{./Images/lith_en_vert_col.pdf}
\maketitle
%\texttt{krikr808@student.liu.se}
%\end{titlepage}

% % % % % % % % % %
% COMPYRIGHT NOTE %
% % % % % % % % % %

\section*{Copyright}
\iffalse
\thispagestyle{empty}
\fi

The publishers will keep this document online on the Internet -- or its possible replacement -- for a period of 25 years starting from the date of publication barring exceptional circumstances.

The online availability of the document implies permanent permission for anyone to read, to download, or to print out single copies for his/hers own use and to use it unchanged for non-commercial research and educational purpose. Subsequent transfers of copyright cannot revoke this permission. All other uses of the document are conditional upon the consent of the copyright owner. The publisher has taken technical and administrative measures to assure authenticity, security and accessibility.

According to intellectual property law the author has the right to be mentioned when his/her work is accessed as described above and to be protected against infringement.

For additional information about the Linköping University Electronic Press and its procedures for publication and for assurance of document integrity, please refer to its www home page: \url{http://www.ep.liu.se/}.
%\hypersetup{colorlinks=true, urlcolor=cyan} % Remove the frame around links

% % % % % % %
% ABSTRACT  %
% % % % % % %

\begin{abstract}
    Your abstract goes here...
\end{abstract}

% Table of contents and various lists
\tableofcontents
\listoftables
\listoffigures

% % % % % % % % %
% INTRODUCTION  %
% % % % % % % % %

\part{Introduction}

\chapter{Motivation}

\chapter{Difficulties}

\begin{itemize}
    \item Wave dispersion
    \item Different depths gives different speeds
    \item The water should interact with moving objects
\end{itemize}

\chapter{Earlier approaches}

\section{Twodimensional methods}

\subsection{Twodimensional PDEs for shallow water}

\subsection{Spectral methods}

References:
\begin{itemize}
    \item \textit{REAL TIME MODELLING OF MULTISPECTRAL OCEAN SCENES}
    \item \textit{GPU-based simulation of Radar sea clutter}
\end{itemize}

\section{Threedimensional methods}

\subsection{Particles}

\subsection{Boundary element method}

\subsection{Finite element method}

\subsection{Finite volume method with tall cells}

\subsection{Finite volume method with octrees}

\section{Twodimensional and threedimensional hybrid methods}

% % % % % %
% METHOD  %
% % % % % %

\part{Method}

\chapter{Finite volume method}

\section{Artificial compression}

\subsection{Instability}

\section{Incompressible Navier-Stokes equations}

\section{Gauss-Seidel method}

\subsection{Preconditioned conjugate gradient method}

\begin{itemize}
    \item Extension of the gradient descent
\end{itemize}

\subsection{Multigrid method}

\section{Semicompressible water}

\chapter{Octrees}

\section{Determining level of detail}

\section{The differentiating problem}

\subsection{Velocity advection term}

\section{Multilevel acceleration}

\section{Wave reflection}

\subsection{Perfectly matched layer}

\chapter{Surface representation}

There are both surface trackibng methods and surface capturing methods.

\section{Level set method}

\subsection{Marching cubes}

\section{Volume of fluid method}

\begin{itemize}
    \item Comparsion: \textit{\href{http://capfluidicslit.mme.pdx.edu/reference/Numerics/Gopala_ChemEngJ2008_VOFMethodsFreeSurfaceFlow.pdf}{Volume of fluid methods for immiscible-fluid and free-surface flows}}
\end{itemize}


\subsection{Interface reconstruction}
%\section{Internal alpha distribution}

\subsection{Geometric advection schemes}

\begin{itemize}
    \item A simple (at least so it seems) scheme: \textit{\href{http://www.lmm.jussieu.fr/~zaleski/nota02.pdf}{A geometrical area-preserving Volume-of-Fluid advection method}}
\end{itemize}

\subsection{Algebraic advection schemes}

\subsubsection{Smearing}

\subsubsection{Convection boundedness criterion (CBC)}

\begin{itemize}
    \item Normalised Variable Diagram (NVD)
    \item \textit{\href{http://warminski.pollub.plwww.ptmts.org.pl/Waclaw-Koron-2-08.pdf}{Comparison of CICSAM and HRIC High-resolution Schemes for Interface Capturing}}
    \item \textit{\href{http://proceedings.fyper.com/eccomascfd2006/documents/85.pdf}{MODELING OF THE WAVE BREAKING WITH CICSAM AND HRIC HIGH-RESOLUTION SCHEMES}}
\end{itemize}

\subsubsection{Hyper-C flux limiter}

\begin{itemize}
    \item Reference: \textit{\href{http://www.water.tkk.fi/wr/kurssit/Yhd-12.112/TVD1.pdf}{The Ultimate Conservative Difference Scheme Applied to Unsteady One-Dimensional Advection}}
\end{itemize}

\paragraph{Floating mixed cells}

\begin{itemize}
    \item Remedy: \textit{\href{https://e-reports-ext.llnl.gov/pdf/245038.pdf}{A Simple Advection Scheme for Material Interface}}
\end{itemize}

\subsubsection{Compressive Interface Capturing Scheme for Arbitrary Meshes (CICSAM)}

\begin{itemize}
    \item Reference: \textit{\href{http://ac.els-cdn.com/S0021999199962769/1-s2.0-S0021999199962769-main.pdf?_tid=85161b57da5f4401e55c9d07495e24ea&acdnat=1336167249_a59e4f578adbacf3bff69936c48cdd57}{A Method for Capturing Sharp Fluid Interfaces on Arbitrary Meshes}}
    \item Test with different Courant numbers: \textit{\href{http://www.marin.nl/upload_mm/8/2/c/1807524470_1999999096_2007-ECCOMAS_HoekstraVazAbeilBunnik.pdf}{Free Surface Flow Modelling with Interface Capturing Techniques}}
    \item Improvement 1: \textit{\href{http://powerlab.fsb.hr/ped/kturbo/openfoam/docs/HenrikRuschePhD2002.pdf}{Computational Fluid Dynamics of Dispersed Two-Phase Flows at High Phase Fractions}}
\end{itemize}

\subsubsection{High Resolution Interface Capturing (HRIC) scheme}

\begin{itemize}
    \item Described here: \textit{\href{http://warminski.pollub.plwww.ptmts.org.pl/Waclaw-Koron-2-08.pdf}{Comparison of CICSAM and HRIC High-resolution Sche\-mes for Interface Capturing}}
\end{itemize}

\subsubsection{Inter-Gamma Scheme}

\begin{itemize}
    \item Reference: \textit{\href{http://powerlab.fsb.hr/ped/kturbo/openfoam/docs/InterTrack.pdf}{Interface Tracking Capabilities of the Inter-Gamma Differencing Scheme}}
\end{itemize}

\subsubsection{Constrained Interpolation Profile (CIP) method}

%TODO: Used for advecting fluid interfaces??

\begin{itemize}
    \item Reference: \textit{\href{http://www.mech.titech.ac.jp/~ryuutai/paper/JCP2001CIPReviewYabe.pdf}{The Constrained Interpolation Profile Method for Multiphase Analysis}}
\end{itemize}

\subsection{Alpha noise due to a non-zero velocity divergence}

\begin{itemize}
    \item Remedy: Advect both water volume and total volume and then define alpha as the ration between them
\end{itemize}

\section{Coupled Level Set/Volume of Fluid  method}

\chapter{Interaction with moving objects}

\chapter{Wind waves}

\section{Spectral methods}

\section{Air-water interaction}

\chapter{Splash and foam}

\chapter{Parallellization}

\section{Space filling curves}

\chapter{Conservation laws}

\begin{itemize}
    \item Mass
    \item Energy
    \item Momentum
    \item Angular momentum
    \item (d/dt)(Center of mass) - momentum = 0
\end{itemize}

% % % % % % % % % % % % % %
% RESULTS AND CONCLUSIONS %
% % % % % % % % % % % % % %

\part{Results and conclusions}

\chapter{Results}

\section{Speed}

\section{Accuracy}

\chapter{Conclusions}

% % % % % % %
% APENDIZES %
% % % % % % %

\part{Appendixes}

\appendix
%\chapter{Derivation of twodimensional PDEs for water waves at random intermediate depths}
\chapter{Derivation of twodimensional PDEs for water waves at varying depths}

% % % % % % % % %
% BIBLIOGRAPHY  %
% % % % % % % % %

\iffalse

\bibliographystyle{plain}
\bibliography{sample1,sample2,...,samplen}
% Note the lack of whitespace between the commas and the next bib file.

\else

\begin{thebibliography}{99}
    
    %TODO: Remove this bibitem
    \bibitem{lamport94}
    Leslie Lamport,
    \emph{\LaTeX: A Document Preparation System}.
    Addison Wesley, Massachusetts,
    2nd Edition,
    1994.
    
\end{thebibliography}

\fi

\end{document}
